\section{The BFS Region-Growing Algorithm}
\label{sec:algorithm}

\subsection{Notation and Subgraph Setup}

Let $G_v = (V_v, E_v)$ be the tract subgraph at bisection-tree node $v$.
Let $p_i$ be tract population and $P_v = \sum_{i \in V_v} p_i$ be node
population. The implemented split is binary: it returns a left set $L$ and a
right set $R$ whose union is $V_v$.

\subsection{Seed Derivation}

\begin{definition}[BFS Node Seed]
\label{def:bfs-node-seed}
For global \texttt{base\_seed} and node path \texttt{path}, BISECT derives
\[
  s_{\mathrm{node}} =
  \mathrm{first8}\left(\mathrm{SHA\text{-}256}
  (\texttt{"BFS\_NODE\_"} \parallel \mathtt{path} \parallel
  \texttt{"\_"} \parallel \mathrm{le8}(\mathtt{base\_seed}))\right).
\]
\end{definition}

\begin{definition}[BFS Split Seed]
\label{def:bfs-split-seed}
Inside \texttt{split\_subgraph\_bfs}, the node seed is domain-separated again:
\[
  s_{\mathrm{split}} =
  \mathrm{first8}\left(\mathrm{SHA\text{-}256}
  (\texttt{"BFS\_SEED\_"} \parallel \mathrm{le8}(s_{\mathrm{node}}))\right).
\]
This seed initializes \texttt{SmallRng} for the population-weighted sample of
the first seed.
\end{definition}

\subsection{Seed Selection}

The split uses exactly two seeds.

\begin{enumerate}
\item \textbf{Seed 0.} Sample a tract from $V_v$ with probability proportional
  to tract population, using the split seed from Definition~\ref{def:bfs-split-seed}.
\item \textbf{Seed 1.} Run BFS from seed 0 and choose a tract with maximum
  finite BFS distance. Ties follow the deterministic iterator order over local
  tract indices.
\end{enumerate}

The result is not a geometric centroid. It is a graph-topological spread: one
seed starts in a population-weighted place, and the other is as far away as the
subgraph permits in hop distance.

\subsection{Priority-Queue Growth}

Each seed owns itself initially. The algorithm pushes unassigned neighbors of
each seed into a binary heap keyed by current district population. Because Rust's
\texttt{BinaryHeap} is a max-heap, the implementation stores
\texttt{Reverse(current\_pop)} so the smaller-population side is popped first.

\begin{definition}[BFS Growth Step]
When a heap entry $(\mathrm{current\_pop}, t, d)$ is popped:
\begin{enumerate}
\item if tract $t$ is already assigned, discard the stale entry;
\item otherwise assign $t$ to side $d$;
\item add $p_t$ to side $d$'s current population; and
\item push every still-unassigned neighbor of $t$ back into the heap with
  side $d$'s updated population.
\end{enumerate}
\end{definition}

The stale-entry rule matters. A tract can be reached from both frontiers; the
first pop wins, later heap entries for the same tract are ignored.

\subsection{Disconnected-Graph Safety Net}

If the induced subgraph is disconnected, the frontier heap can empty before
every tract is assigned. The implementation then assigns any leftover tract to
the side with smaller current population. This makes the output complete, but
it weakens any pure contiguity guarantee. Therefore the correct claim is:
\bfsg{} is contiguous by construction on connected subgraphs before the
disconnected-graph safety net and subject to any later rebalance moves.

\subsection{Post-Hoc Rebalance}

After growth, the splitter performs up to 200 local boundary moves. At each
iteration it computes left population, compares it to $P_v/2$, identifies the
heavy side, and moves one boundary tract whose population best reduces the
excess if such a boundary tract exists. The current implementation does not run
a full contiguity check inside this rebalance loop; it uses boundary adjacency
as the local repair signal.

This is enough for a fast construction baseline, but not enough for a strong
legal-style contiguity theorem. The final plan should still pass the platform's
ordinary validation surface.

\subsection{Formal Pseudocode}

\begin{algorithm}[t]
\caption{\bfsg{}-\textsc{Bisect}($G_v$, $\mathbf{p}$, $\mathtt{base\_seed}$, $\mathtt{path}$)}
\label{alg:bfs}
\begin{algorithmic}[1]
\State $s_{\mathrm{node}} \gets \mathrm{BFSNodeSeed}(\mathtt{base\_seed}, \mathtt{path})$
\State $s_{\mathrm{split}} \gets \mathrm{BFSSplitSeed}(s_{\mathrm{node}})$
\State $\mathrm{seed}_0 \gets \mathrm{WeightedSample}(V_v, \mathbf{p}, s_{\mathrm{split}})$
\State $\mathrm{seed}_1 \gets \arg\max_t d_{\mathrm{BFS}}(\mathrm{seed}_0,t)$
\State assign seed 0 to side 0 and seed 1 to side 1
\State push their unassigned neighbors into a heap keyed by side population
\While{heap is non-empty}
  \State pop the entry whose side has smaller current population
  \If{tract is already assigned} \textbf{continue} \EndIf
  \State assign tract to that side and update side population
  \State push unassigned neighbors for the same side
\EndWhile
\State assign disconnected leftovers to the currently smaller side
\State apply up to 200 boundary rebalance moves toward $P_v/2$
\State \textbf{return} left and right tract sets
\end{algorithmic}
\end{algorithm}

\subsection{Runtime}

On a connected node, each edge can generate heap activity from frontier growth,
so the natural bound is $O((|V_v|+|E_v|)\log |V_v|)$ plus the bounded 200-step
rebalance pass. On planar tract graphs this behaves like $O(n\log n)$ per
bisection node. End-to-end statewide runtime depends on the bisection tree and
validation pipeline.
